\section{Excluded-coordinate translations on a finite circle}
\label{sec:excluded-coordinate}\label{sec:excluded}

We now construct a second family of hard binary words.  The parameter is a
sufficiently large integer $m$; the eventual distortion is at least
$m^{m/4}$, while the logarithm of the word length is $O(m^2\log m)$.
At each level there is one prime for each of the $m$ child indices.  The
translation associated with that prime acts on every child except its own.
This omission has two
uses.  A small phase cost will force most translation coordinates to vanish,
whereas a frequency annihilating all the translations and detecting the
half-turn must have large total absolute leaf frequency. The proof of this
frequency bound assumes the total is small and forces branching through
too many children.  We first prove the phase assertion, then realize it
by words and compare it with the frequency assertion.

\subsection{The translation pattern and its phase cost}

Write $\mathbb T=\mathbb R/\mathbb Z$ for the circle of circumference one,
and let $\|z\|_{\mathbb T}$ denote the distance of $z$ from zero.  Set $h=m$.
We need $hm=m^2$ primes, all distinct across all levels, satisfying
\begin{equation}\label{eq:exc-primes}
 m^3<p_{\ell,i}\le m^5,
 \qquad 1\le\ell\le h,\quad 1\le i\le m.
\end{equation}
The following count uses the central binomial coefficient and Legendre's
factorial valuation formula in the manner of Erd\H{o}s \cite{Erdos1932},
who attributes the underlying ideas to Landau.  The central term is the
largest of the $2n+1$ binomial coefficients summing to $4^n$, so
\[
 \binom{2n}{n}\ge\frac{4^n}{2n+1}.
\]
In the valuation formula for $\binom{2n}{n}$, each term
$\lfloor2n/p^j\rfloor-2\lfloor n/p^j\rfloor$ is zero or one.  Thus the
exponent of $p$ is at most $\lfloor\log_p(2n)\rfloor$, and each prime
contributes a factor at most $2n$.  Writing $\pi(x)$ for the number of
primes at most $x$, we obtain
\[
 \binom{2n}{n}\le (2n)^{\pi(2n)}.
\]
Taking logarithms and then $n=\lfloor m^5/2\rfloor$ gives
$\pi(m^5)\ge c m^5/\log m$ for an absolute $c>0$ and all sufficiently
large $m$.  Since $\pi(m^3)\le m^3$, the interval $(m^3,m^5]$ contains
at least $m^2$ primes once $m$ is large enough.  Fix such a choice.

For each level put
\begin{equation}\label{eq:exc-translations}
 P_\ell=\prod_{i=1}^m p_{\ell,i},\qquad
 u_{\ell,i}(s)=\frac{s}{p_{\ell,i}}\pmod1,\qquad
 c_{\ell,i}(s)=\sum_{j\ne i}u_{\ell,j}(s).
\end{equation}
The parameter $s$ is taken modulo $P_\ell$.  By the Chinese remainder
theorem, $(u_{\ell,i}(s))_{i=1}^m$ runs through the full product of the
prime-order subgroups.  In particular, one coordinate can be prescribed
while all the others are set to zero.  Each $c_{\ell,i}$ is a homomorphism.
If only $u_{\ell,i}$ is nonzero, its value is added to all children except
child $i$.

The scalar cost of a phase after $\ell$ levels is
\begin{equation}\label{eq:exc-rho}
 \rho_0(z)=\|z\|_{\mathbb T},\qquad
 \rho_\ell(z)=\min_{s\bmod P_\ell}
       \frac1m\sum_{i=1}^m
          \rho_{\ell-1}\bigl(z+c_{\ell,i}(s)\bigr)
       \quad(1\le\ell\le h).
\end{equation}
These functions are defined on all of $\mathbb T$.  Each minimum is over a
finite set.
They are symmetric and subadditive: negate a minimizing parameter for
symmetry, and add minimizing parameters for two inputs and use the preceding
subadditivity term by term.  The zero parameter shows inductively that, for
$1\le\ell\le h$,
\[
 \rho_\ell(0)=0,\qquad
 0\le \rho_\ell(z)\le \rho_{\ell-1}(z)
       \le\|z\|_{\mathbb T}\le\frac12.
\]

The next lemma quantifies what the omitted coordinate prevents.  Its input
phase has a small order relatively prime to the primes still available in
the recurrence.  The phases followed during the proof can have much larger
orders.

\begin{lemma}[Small-order phases remain costly]\label{lem:exc-small-cost}
For $0\le\ell\le h$, put
\[
 F_\ell=m^{-\ell/4-100m^{3/4}}.
\]
For all sufficiently large $m$, the following holds for every
$0\le\ell\le h$.  If $z\in\mathbb T$ is nonzero, has finite order
$q\le m^{10}$, and satisfies
$\gcd(q,\prod_{a=1}^{\ell}P_a)=1$, then
\begin{equation}\label{eq:exc-phase-lower}
 \rho_\ell(z)\ge F_\ell.
\end{equation}
The empty product in the case $\ell=0$ is one.
\end{lemma}

\begin{proof}
At height zero a nonzero point of order $q$ has distance at least
$1/q\ge m^{-10}\ge F_0$ for large $m$.  Suppose the lemma is known at all
heights below $\ell\ge1$, and suppose, for a contradiction, that an input
as in the statement satisfies $\rho_\ell(z)<F_\ell$.

We record the bounds that will keep a descent through the recurrence under
control.  Let $a=m^{-1/4}$ and $b_t=4a^{t+1}$ for $t\ge0$.  For
$m\ge32^4$, all $b_t<1$ and
\begin{align}
 \log\prod_{t\ge0}(1-b_t)^{-1}
 &\le \sum_{t\ge0}\frac{b_t}{1-b_t}
 \le \frac{4a}{(1-a)(1-4a)}<\log2,
 \label{eq:exc-product-cost}\\
 m\sum_{t\ge0}b_t
 &=\frac{4m^{3/4}}{1-m^{-1/4}}\le5m^{3/4}.
 \label{eq:exc-support-budget}
\end{align}
We will choose phases $z_0=z,z_1,\ldots,z_\ell$ so that, at time $t$,
\begin{equation}\label{eq:exc-path-invariant}
 \rho_{\ell-t}(z_t)
 <F_\ell\prod_{j<t}(1-b_j)^{-1}<2F_\ell .
\end{equation}
The first inequality holds at $t=0$ by assumption.

Suppose $t<\ell$ and write $r=\ell-t$ for the remaining height.  Choose a
parameter attaining the minimum in $\rho_r(z_t)$, abbreviate
$u_i=u_{r,i}(s)$, and put $U=\sum_i u_i$.  Call an index $i$ small when
\[
 \rho_{r-1}(z_t+U-u_i)<F_{r-1}/2.
\]
The average of these $m$ costs is less than $2F_\ell$ by
\eqref{eq:exc-path-invariant}.  Hence the number of indices that are not
small is at most
\begin{equation}\label{eq:exc-exception-count}
 \frac{4mF_\ell}{F_{r-1}}
   =4m^{1-(t+1)/4}=mb_t.
\end{equation}
This is at most $4m^{3/4}$, so there are at least two small indices for all
large $m$.

If two small indices $i,j$ had $u_i\ne u_j$, symmetry and subadditivity
would give
\[
 \rho_{r-1}(u_i-u_j)
 \le \rho_{r-1}(z_t+U-u_i)+\rho_{r-1}(z_t+U-u_j)
 <F_{r-1}.
\]
The difference is nonzero and its order divides
$p_{r,i}p_{r,j}\le m^{10}$.  That order is relatively prime to
$\prod_{a<r}P_a$, since all level primes are distinct.  The induction
hypothesis at height $r-1$ rules out the displayed inequality.  Therefore
all small indices have the same $u_i$.  Two distinct prime-order subgroups
of $\mathbb T$ intersect only at zero; the presence of two small indices
then implies that every small index has $u_i=0$.

Let $k_t=\#\{i:u_i\ne0\}$.  We have $k_t\le mb_t<m$.  Every zero coordinate
gives the same child phase $z_t+U$; choose
$z_{t+1}=z_t+U$.  The other costs are nonnegative, so
\begin{equation}\label{eq:exc-path-cost}
 \rho_{r-1}(z_{t+1})
 \le\frac{\rho_r(z_t)}{1-k_t/m}
 \le\frac{\rho_r(z_t)}{1-b_t}.
\end{equation}
This proves the first inequality of \eqref{eq:exc-path-invariant} at
$t+1$; its second inequality follows from \eqref{eq:exc-product-cost}.
Thus the descent is justified at every step.  It uses fractions from at
most
\begin{equation}\label{eq:exc-denominator-count}
 \sum_{t=0}^{\ell-1}k_t
 \le m\sum_{t\ge0}b_t\le5m^{3/4}
\end{equation}
distinct level primes, by \eqref{eq:exc-support-budget}.

It remains to compare the terminal phase with its cost.  Let $Q$ be the
product of the primes actually used in the descent.  Then $z_\ell-z$ has
order dividing $Q$, and
\[
 qQ\le m^{10+25m^{3/4}}.
\]
The phase $z_\ell$ is nonzero.  Otherwise the order $q>1$ of $z$ would
divide $Q$, contradicting $\gcd(q,Q)=1$.  Since the order of $z_\ell$
divides $qQ$, we conclude that
\[
 \rho_0(z_\ell)=\|z_\ell\|_{\mathbb T}
 \ge m^{-10-25m^{3/4}}.
\]
But \eqref{eq:exc-path-invariant} gives
\[
 \rho_0(z_\ell)<2F_\ell
   =2m^{-\ell/4-100m^{3/4}}
   <m^{-10-25m^{3/4}}
\]
for sufficiently large $m$, a contradiction.  The induction hypothesis
was used only for the differences $u_i-u_j$.  The selected phases $z_t$
were controlled by their accumulated denominators, without a small-order
assumption on them.
\end{proof}

\subsection{Words whose tags preserve the child index}

We now turn the scalar cost into edit separation.  Put
\begin{equation}\label{eq:exc-size-parameters}
 L=2\prod_{\ell=1}^hP_\ell,\qquad R=m^{2h},\qquad
 \Lambda=[m]^h,\qquad
 G=(L^{-1}\mathbb Z/\mathbb Z)^\Lambda .
\end{equation}
The set $\Lambda$ indexes the leaves of an $m$-ary tree of height $h$:
a node at height $\ell$ is a prefix $a\in[m]^{h-\ell}$, and $ai$ denotes
its child obtained by appending $i$.  The alphabet is
\[
 \mathcal A=\Lambda\times\{0,\ldots,L-1\}\times\{0,1\},
 \qquad J=|\mathcal A|=2Lm^h.
\]
The first two coordinates are symbol tags.  The factor two in $L$ places
the half-turn in the phase group, and every translation
$c_{\ell,i}(s)$ also belongs to this group.

For a state $v\in G$ and an auxiliary phase
$\theta\in L^{-1}\mathbb Z/\mathbb Z$, define $W_a(v,\theta)$ recursively.
At a leaf $a$, it is the length-$L$ word
\begin{equation}\label{eq:exc-leaf-word}
 W_a(v,\theta)=
 \bigl(a,b,\mathbf1_{[0,1/2)}(b/L-v_a-\theta)\bigr)_{b=0}^{L-1},
\end{equation}
where the half-open interval is interpreted on $\mathbb T$.  At height
$\ell\ge1$, concatenate the following rows, each consisting of $m$ child
words:
\begin{equation}\label{eq:exc-internal-word}
 W_a(v,\theta)=
 \mathop{\Vert}_{s=0}^{RP_\ell-1}
 \mathop{\Vert}_{i=1}^{m}
 W_{ai}\bigl(v,\theta+c_{\ell,i}(s)\bigr).
\end{equation}
The symbol $\Vert$ denotes concatenation in the displayed order.  All
height-$\ell$ words have the same length
\begin{equation}\label{eq:exc-length-recurrence}
 N_0=L,\qquad N_\ell=mRP_\ell N_{\ell-1}.
\end{equation}
Write $N=N_h$ and $W(v)=W_{\varnothing}(v,0)$.  Adding a phase $z$ to
every state coordinate below a node has the same effect on its word as
replacing $\theta$ by $\theta+z$: this follows at a leaf from
\eqref{eq:exc-leaf-word}, and then at every height from the concatenation.

\begin{lemma}[Separation under arbitrary alignments]\label{lem:exc-alignment}
For every node $a$ at height $0\le\ell\le h$, every $v\in G$, and every
$\theta,z\in L^{-1}\mathbb Z/\mathbb Z$,
\begin{equation}\label{eq:exc-word-separation}
 r_{N_\ell}\bigl(W_a(v,\theta),W_a(v,\theta+z)\bigr)
 \ge N_\ell 3^{-\ell}\rho_\ell(z).
\end{equation}
\end{lemma}

\begin{proof}
At a leaf, the position tag $b$ forces a matched symbol to use the same
position on both sides.  The two half-circle indicators disagree at
exactly $2L\|z\|_{\mathbb T}$ positions.  Their deficit is therefore
exactly that number, which is at least $L\rho_0(z)$.

Suppose the lemma holds at height $\ell-1$ and fix a longest common
subsequence at height $\ell$, with deficit $E$.  A row in
\eqref{eq:exc-internal-word} has $m$ children of length $N_{\ell-1}$.
Because every symbol carries its leaf tag, a match between these two node
words preserves the child index $i$.  We verify the nonsplit-row hypothesis
of Lemma~\ref{pre:tagged-rows}.

Suppose all matches from a source row with parameter $s$ lie in a target
row with parameter $s'$.  Within child $i$ they form a matching of two full
child words whose phase difference is
$z+c_{\ell,i}(s'-s)$.  By induction, the number of unmatched source
positions in that child is at least
\[
 N_{\ell-1}3^{-(\ell-1)}
       \rho_{\ell-1}\bigl(z+c_{\ell,i}(s'-s)\bigr).
\]
Summing over $i$ and using \eqref{eq:exc-rho}, the source row loses at least
the fraction
\[
 A=3^{-(\ell-1)}\rho_\ell(z)
\]
of its positions.  A row with no matches loses every position and satisfies
the same estimate.  Here $0\le A\le1/2$.  Lemma~\ref{pre:tagged-rows},
with $t=m$ and child length $N_{\ell-1}$, now gives
\[
 E\ge\frac{A N_\ell}{1+mA/(m-1)}
   \ge \frac{A N_\ell}{3}.
\]
For the last inequality we used $m/(m-1)\le2$ and $A\le1$.  This is
\eqref{eq:exc-word-separation}.  Lemma~\ref{pre:tagged-rows} applies to the
chosen arbitrary matching; in particular it permits different nonsplit
source rows to use the same target row.
\end{proof}

Let $w_*\in G$ have every coordinate equal to $1/2$.  All primes in
\eqref{eq:exc-primes} are odd.  The phase $1/2$ has order two, so
Lemmas~\ref{lem:exc-small-cost} and~\ref{lem:exc-alignment}, together with
the phase identity for node words, give the pointwise bound
\begin{equation}\label{eq:exc-halfturn}
 r_N\bigl(W(v),W(v+w_*)\bigr)\ge N3^{-h}F_h
 \qquad(v\in G).
\end{equation}

\subsection{Cheap translations and the binary transfer}

For a node $a$ at height $\ell\ge1$ and $s\bmod P_\ell$, define
$\tau_{a,s}\in G$ by
\[
 (\tau_{a,s})_b=
 \begin{cases}
 c_{\ell,i}(s),&\text{if the leaf $b$ has prefix $ai$},\\
 0,&\text{if $b$ is outside the subtree below $a$}.
 \end{cases}
\]
Let $H\le G$ be the subgroup generated by these vectors over all internal
nodes.  At each fixed node, $s\mapsto\tau_{a,s}$ is a homomorphism.

Fix one height $\ell$ and choose one translation at each node of that
height.  In an occurrence of $W_a(v,\theta)$, adding $\tau_{a,s}$ replaces
row parameter $t$ by $t+s$ in \eqref{eq:exc-internal-word}.  The row list
has period $P_\ell$ and repeats it $R$ times.  Choose $0\le s<P_\ell$.
Deleting the first $s$ rows from the unshifted list and the last $s$ rows
from the shifted list leaves identical words.  The edit cost in this
occurrence is at most
\[
 2s\,mN_{\ell-1}\le \frac{2N_\ell}{R}.
\]
At a fixed height the node occurrences partition the root word.  The
combined cost for all nodes at that height is therefore at most $2N/R$.
Every element of $H$ is a sum of at most $h$ such single-height
translations: $G$ is abelian, sums at one node combine to one parameter,
and the nodes can then be grouped by height.  The edit estimate holds for
every state and every ancestor phase.  Applying it at each intermediate
state and using the triangle inequality gives
\begin{equation}\label{eq:exc-cheap-subgroup}
 \operatorname{ED}\bigl(W(v),W(v+u)\bigr)
 \le\frac{2hN}{R}
 \qquad(v\in G,\ u\in H).
\end{equation}

For a leaf $a\in\Lambda$, write $ze_a$ for the state vector with value $z$
at $a$ and zero elsewhere. This shift changes exactly $2L\|z\|_{\mathbb T}$ symbols in every occurrence
of that leaf word.  Each branching has $m$ equally long children, so the
occurrences of a fixed leaf occupy the fraction $m^{-h}$ of the root word.
Substituting only the changed symbols yields
\begin{equation}\label{eq:exc-cheap-coordinate}
 \operatorname{ED}\bigl(W(v),W(v+ze_a)\bigr)
 \le\frac{2N}{m^h}\|z\|_{\mathbb T}.
\end{equation}
This holds for every $v\in G$, $z\in L^{-1}\mathbb Z/\mathbb Z$, and
$a\in\Lambda$.

Apply the second construction of Theorem~\ref{code:two-proofs}, proved
using anchors at equal relative indices in Section~\ref{p2:sec:code-anchor}, to
$\mathcal A$.  It gives one code
$\mathcal C:\mathcal A\to\{0,1\}^{w}$ with
$w\le C_{\rm code}\log(2J)$ and an absolute $\eta>0$.  Its exact estimates
are
\begin{equation}\label{eq:exc-code-interface}
 \begin{aligned}
 r_{wq}(\mathcal Cx,\mathcal Cy)&\ge\eta w r_q(x,y)
       &&(x,y\in\mathcal A^q),\\
 \operatorname{ED}(\mathcal Cx,\mathcal Cy)
       &\le w\operatorname{ED}(x,y).
 \end{aligned}
\end{equation}
The lower estimate holds simultaneously at every common length $q\ge0$,
with the code fixed before $q$ is chosen; the upper estimate holds for
arbitrary finite words.  Here we use the lower estimate at the common
length $N$.

Set $X(v)=\mathcal C W(v)$ and $n_*=wN$.  The metric comparison
$r_{n_*}\le\operatorname{ED}$ from Lemma~\ref{pre:metric}, the code
estimates, and \eqref{eq:exc-halfturn}--\eqref{eq:exc-cheap-coordinate}
give
\begin{align}
 \frac{\operatorname{ED}(X(v),X(v+w_*))}{n_*}
 &\ge\eta3^{-h}F_h,
 \label{eq:exc-binary-far}\\
 \frac{\operatorname{ED}(X(v),X(v+u))}{n_*}
 &\le\frac{2h}{R},
 \label{eq:exc-binary-subgroup}\\
 \frac{\operatorname{ED}(X(v),X(v+ze_a))}{n_*}
 &\le\frac{2}{m^h}\|z\|_{\mathbb T}.
 \label{eq:exc-binary-coordinate}
\end{align}
All three inequalities hold for every $v\in G$.  The second holds for every
$u\in H$, and the third for every $z\in L^{-1}\mathbb Z/\mathbb Z$ and
$a\in\Lambda$.  In particular, each half-turn pair consists of distinct
words; other repetitions among the state images are allowed.

\subsection{Large total leaf frequency}

It remains to show that the cheap translations prevent an $\ell_1$ map
from preserving the half-turn distance.  Write characters of $G$ as
\[
 \chi_\xi(v)=\exp\!\left(2\pi\mathrm i\sum_{a\in\Lambda}\xi_a v_a\right),
 \qquad
 \xi\in(\mathbb Z/L\mathbb Z)^\Lambda .
\]
For each coordinate choose an integer representative of minimum absolute
value, with either sign allowed in a tie, and put
\begin{equation}\label{eq:exc-frequency-size}
 T(\xi)=\sum_{a\in\Lambda}|\xi_a|.
\end{equation}
Thus $T$ records the total magnitude of the leaf frequencies.  Let
$H^\perp=\{\xi:\chi_\xi(u)=1\text{ for every }u\in H\}$.

\begin{lemma}[Annihilating frequencies that detect the half-turn]
\label{lem:exc-annihilator}
If $\xi\in H^\perp$ and $\chi_\xi(w_*)\ne1$, then
\begin{equation}\label{eq:exc-annihilator-threshold}
 T(\xi)\ge K,\qquad K=m^{3h/4}.
\end{equation}
\end{lemma}

\begin{proof}
For each node $a$, let $t_a$ be the sum of the chosen integer
representatives over all leaves below $a$.  Thus $t_a=\xi_a$ at a leaf,
and $t_a=\sum_i t_{ai}$ at an internal node.  At height $\ell$, the Chinese
remainder theorem permits a parameter with
$u_{\ell,i}=1/p_{\ell,i}$ and all other $u_{\ell,j}=0$.  Its node
translation adds $1/p_{\ell,i}$ to every child other than $i$.  Since
$\xi$ annihilates this translation,
\begin{equation}\label{eq:exc-annihilator-congruence}
 t_a-t_{ai}\equiv0\pmod{p_{\ell,i}}
 \qquad(1\le i\le m).
\end{equation}
This congruence is independent of the integer lifts because every
$p_{\ell,i}$ divides $L$.

The half-turn satisfies $\chi_\xi(w_*)=(-1)^{t_{\varnothing}}$.
Consequently the root sum is odd and is nonzero.  Suppose that $T(\xi)<K$.
Every node sum has absolute value at most $T(\xi)$, hence less than $K$.
At a node with $t_a\ne0$, let $q$ be the number of its level primes
dividing $t_a$.  If $q>0$, distinctness and \eqref{eq:exc-primes} give
\[
 m^{3q}
 <\prod_{p_{\ell,i}\mid t_a}p_{\ell,i}
 \le |t_a|<m^{3h/4}.
\]
Thus $q<h/4$; the same inequality is immediate when $q=0$.
For every other child, \eqref{eq:exc-annihilator-congruence} forces
$t_{ai}\ne0$.  Since $h=m$, each nonzero internal node has at least
$3m/4$ nonzero children.  Starting with the root, this gives at least
$(3m/4)^h$ nonzero leaf coefficients.  These are nonzero integers, so
\[
 T(\xi)\ge(3m/4)^h\ge m^{3h/4}=K
\]
for all sufficiently large $m$, a contradiction.
\end{proof}

\subsection{The distortion and length estimates}

Let $f$ be any injection of the finite set $X(G)$ into $\ell_1$, with
distortion $D$.  Rescale it by its contraction factor so that
\begin{equation}\label{eq:exc-normalization}
 \operatorname{ED}(x,y)\le\|f(x)-f(y)\|_1
       \le D\operatorname{ED}(x,y)
 \qquad(x,y\in X(G)).
\end{equation}
This applies to each injection individually.  Put $g(v)=f(X(v))/n_*$
and, using uniform measure on $G$, define
\[
 \mathcal E(u)=\mathbb E_{v\in G}\|g(v+u)-g(v)\|_1 .
\]
The cut--Fourier representation in Lemma~\ref{pre:weighted-circle}
provides nonnegative weights $\omega_\xi$ such that
\begin{equation}\label{eq:exc-fourier-energy}
 \mathcal E(u)=\sum_\xi\omega_\xi|1-\chi_\xi(u)|^2 .
\end{equation}
The lemma applies with period $L$, $m^h$ coordinates, subgroup $H$,
and threshold $K$ by Lemma~\ref{lem:exc-annihilator}.  Its representation
and intermediate weight estimates use normalized counting measure and
also apply when the target $\ell_1$ has infinitely many coordinates.

We spell out the two resulting bounds on frequency mass.  Let
$\mathcal D_L=\{1,2,4,\ldots\}\cap[1,L/2]$.  For every $s\in\mathcal D_L$,
\eqref{eq:exc-normalization} and \eqref{eq:exc-binary-coordinate} imply
\[
 \frac Ls\sum_{a\in\Lambda}\mathcal E((s/L)e_a)
 \le \frac Ls\,m^h\frac{2D}{m^h}\frac{s}{L}=2D.
\]
The weighted estimate in Lemma~\ref{pre:weighted-circle} therefore
gives
\begin{equation}\label{eq:exc-weighted-mass}
 \begin{aligned}
 4\sum_\xi\omega_\xi T(\xi)
 &\le\sum_{s\in\mathcal D_L}\frac Ls
           \sum_{a\in\Lambda}\mathcal E((s/L)e_a)\\
 &\le 2D|\mathcal D_L|
 \le C D\log(L+1).
 \end{aligned}
\end{equation}
In particular $\sum_\xi\omega_\xi T(\xi)\le C D\log(L+1)$ after adjusting
the absolute constant.  The subgroup identity in that lemma, together
with \eqref{eq:exc-binary-subgroup}, gives
\begin{equation}\label{eq:exc-off-annihilator}
 \sum_{\xi\notin H^\perp}\omega_\xi
 =\frac12\mathbb E_{u\in H}\mathcal E(u)
 \le \frac{Dh}{R}.
\end{equation}
This is the uniform average on the finite subgroup $H$ itself.

The half-turn multiplier in \eqref{eq:exc-fourier-energy} equals four
when $t_{\varnothing}$ is odd and zero otherwise.  A detecting frequency
inside $H^\perp$ has $T(\xi)\ge K$, whereas
\eqref{eq:exc-off-annihilator} controls all frequencies outside
$H^\perp$.  Using also the noncontracting half-turn bound from
\eqref{eq:exc-binary-far}, we obtain
\begin{align}
 \eta3^{-h}F_h
 &\le \mathcal E(w_*) \notag\\
 &\le 4\sum_{\xi\notin H^\perp}\omega_\xi
       +\frac4K\sum_\xi\omega_\xi T(\xi) \notag\\
 &\le C D\left(\frac hR+\frac{\log(L+1)}K\right).
 \label{eq:exc-final-energy}
\end{align}
Thus the separation in the word metric and the branching of frequencies
have now been compared using the same finite subgroup.

The prime bounds give
\begin{equation}\label{eq:exc-logL}
 \log L\le\log2+5hm\log m.
\end{equation}
With $h=m$, $R=m^{2h}$, and $K=m^{3h/4}$, the term $h/R$ in
\eqref{eq:exc-final-energy} is bounded by $\log(L+1)/K$ for all large $m$.
Since $F_h=m^{-h/4-100m^{3/4}}$, it follows that
\begin{equation}\label{eq:exc-distortion-explicit}
 D\ge c\eta\,
       \frac{3^{-m}m^{m/2-100m^{3/4}}}{\log(L+1)}
\end{equation}
for an absolute $c>0$.  The logarithm of the right side, divided by
$m\log m$, tends to $1/2$: the losses contribute
$100m^{-1/4}$, $\log3/\log m$, and a term tending to zero by
\eqref{eq:exc-logL}.  Therefore, after increasing an absolute threshold,
\begin{equation}\label{eq:exc-distortion-power}
 D\ge m^{m/4}\qquad(m\ge m_0).
\end{equation}
Because $f$ was arbitrary, this is a distortion lower bound for the finite
set $X(G)$.

Finally, \eqref{eq:exc-size-parameters} and
\eqref{eq:exc-length-recurrence} give the exact root length
\[
 N=L\prod_{\ell=1}^h(mRP_\ell)
   =\frac{L^2}{2}(mR)^h.
\]
Consequently
\[
 \log N\le12m^2\log m+m\log m+\log2.
\]
Also $J=2Lm^h$ and $w\le C_{\rm code}\log(2J)$, so
$\log w=O(\log m)$.  Hence, for an absolute $C_0$,
\begin{equation}\label{eq:exc-binary-length}
 \log n_*=\log(wN)\le C_0m^2\log m
 \qquad(m\ge m_0).
\end{equation}
We have obtained these binary same-length witnesses for every sufficiently
large integer $m$.  Section~\ref{sec:assembly} applies its direct
all-length inversion to \eqref{eq:exc-distortion-power} and
\eqref{eq:exc-binary-length}, yielding a witness of one length at most
each sufficiently large cap.
